\section{Discussion}
\label{sec:discussion}

\subsection{Revisiting the Expected Findings}
Our project proposal stated six expectations. We revisit them in the light of the results.

\begin{enumerate}
  \item \textbf{Multimodal models outperform unimodal ones on at least some tasks.} \emph{Supported on the larger
  datasets.} Fusion significantly beats the best text model on all five BAM and MemoSen tasks and the best image
  model on two of them. On the 133 labelled CMBan memes the evidence is inconclusive. The gain over the image model
  is small (1.5--2 points), because CLIP already captures much of what the text adds.

  \item \textbf{Humour and sarcasm remain hard.} \emph{Partly supported.} Sentiment, not humour or sarcasm, is the
  hardest task on every dataset. However, CMBan Humour and Sarcasm have the largest share of contextual/cultural
  errors after Sentiment (about 20\%), and BAM Sarcasm is the task where fusion gains least over the image model
  (+0.2 points).

  \item \textbf{Different perturbations affect different models and tasks differently.} \emph{Supported.}
  Linguistic perturbations hurt text and fusion heads, visual ones hurt image and fusion heads, and the effects
  differ by task. Romanisation costs CMBan Offensive 13--18 F1 points and BAM Vulgarity 10--13 points, but leaves
  Sarcasm almost unchanged on both datasets. Tasks signalled by specific words (swear words, slurs) are the most
  vulnerable to a change of script. Sarcasm depends more on the incongruity between image and text, and is the least
  vulnerable.

  \item \textbf{Linguistic perturbations hurt text-heavy memes more; visual corruption hurts where the image is the
  evidence.} \emph{Supported.} Romanisation hurts Bangla-script memes most, whose entire text is rewritten. Visual
  corruptions matter most on the datasets where the image model is strongest (BAM, MemoSen).

  \item \textbf{Explainability reveals modality-specific failure modes.} \emph{Supported, with a surprise.} IG shows
  that Romanised text attracts \emph{more} attribution rather than less. Grad-CAM shows over-attention to rendered
  captions on cartoon memes, but not on photographic memes.

  \item \textbf{The most vulnerable perturbations and dimensions can be identified.} \emph{Supported.} Romanisation
  is the dominant threat on every dataset, followed by blur and down-sampling on the larger datasets. Lexically signalled
  dimensions (offensiveness, vulgarity) and sentiment are the most affected.
\end{enumerate}

\subsection{Why Romanisation Breaks Fusion Models}
Three observations combine into a consistent explanation.
\begin{enumerate}
  \item XLM-R's representation of Romanised Bangla is far from that of Bangla script. Native-script spelling noise
  of similar strength is harmless.
  \item The fusion head has never seen such inputs during training. It treats the out-of-distribution text vector
  as strong evidence, which is reflected in the increased IG text share.
  \item With little training data (CMBan), the head has not learned to rely on the image, so the error propagates
  fully. With more data, the image branch absorbs about half of the damage.
\end{enumerate}
The augmentation experiment (Section~\ref{sec:results:aug}) confirms this account. Showing the head transliterated
copies of its training memes removes most of the drop, including for an unseen Romanisation style. The weakness
therefore lies in the training distribution of the head, not in an inability of the encoder to represent Banglish.

\subsection{Practical Recommendations}
\begin{itemize}
  \item \textbf{Augment with Romanised text.} Transliteration augmentation is cheap, needs no extra labels, and
  removes most of the dominant failure mode at a cost of under one F1 point on clean data.
  \item \textbf{Report flip rates.} Visual perturbations change up to 27\% of predictions with only a few points of
  F1 loss. Aggregate accuracy alone overstates stability.
  \item \textbf{Do not discard character features.} Character $n$-gram models beat frozen sentence encoders on
  meme text. Combining both, or fine-tuning the text encoder, is a promising direction.
  \item \textbf{Prefer a stronger image encoder when affordable.} CLIP ViT-L/14 gives the largest single accuracy
  gain (2.7--2.8 points). The image branch is also what cushions text corruption.
\end{itemize}

\subsection{CMBan in Context}
CMBan's cartoon memes differ from the photographic memes of BAM and MemoSen in three ways:
\begin{itemize}
  \item Their text is longer (24.8 vs.\ about 15 words).
  \item CLIP attends more to their captions.
  \item Their errors are more often contextual or cultural (18\% vs.\ 10--11\%).
\end{itemize}
These differences support CMBan's motivation as a distinct resource: cartoon memes rely heavily on the written
dialogue and on shared cultural knowledge of recurring templates. They also mean that conclusions drawn on
photographic Bangla memes do not automatically transfer to cartoon memes. The full CMBan label release would allow
these differences to be tested with adequate statistical power.
